\chapter{Tiny210 根文件系统、bcache
与日志}\label{tiny210-ux6839ux6587ux4ef6ux7cfbux7edfbcache-ux4e0eux65e5ux5fd7}

\section{当前根文件系统来源}\label{ux5f53ux524dux6839ux6587ux4ef6ux7cfbux7edfux6765ux6e90}

Tiny210 的根文件系统格式仍是 xv6fs，但后端按以下优先级选择：

\begin{lstlisting}
SD 卡原始 xv6 分区
  -> FAT32 启动分区中的 FS.IMG
  -> 都不可用时，回退到内核 .ramdisk 中的内嵌 fs.img
\end{lstlisting}

实板构建已经编入 SDHCI、FAT32、VFS 和只读 ext2。FAT32 启动分区可挂载到
\passthrough{\lstinline!/boot!}；兼容 ext2 分区可通过 VFS 挂载读取。QEMU
替身构建没有 S5PV210 SDHCI， 仍使用内嵌 RAM 镜像。

\section{镜像生成}\label{ux955cux50cfux751fux6210}

\href{../../../ports/tiny210/Makefile}{\passthrough{\lstinline!ports/tiny210/Makefile!}}
先编译 ARM32 用户程序，再调用
\href{../../../mkfs/mkfs.c}{\passthrough{\lstinline!mkfs/mkfs.c!}}：

\begin{lstlisting}
UPROGS
  -> _init, _armcheck, _sh, _echo, _cat, _ls, ...
  -> mkfs fs.img UPROGS...
       创建根 inode 1
       创建 /bin inode
       所有命令放到 /bin
       /init 与 /bin/init 是同一 inode 的两个硬链接
  -> .incbin "fs.img"
  -> kernel .ramdisk
\end{lstlisting}

默认 \passthrough{\lstinline!TINY210\_FSBLOCKS=4096!}，块大小
\passthrough{\lstinline!BSIZE=1024!}，总容量 4 MiB。可在构建时调整：

\begin{lstlisting}[language=sh]
make -f ports/tiny210/Makefile TINY210_FSBLOCKS=8192
\end{lstlisting}

增大文件系统也会直接增大内核/uImage 和 BL1 从 SD
读取的扇区数，必须确保镜像仍 位于预留原始区内。

\section{磁盘格式}\label{ux78c1ux76d8ux683cux5f0f}

xv6fs 使用固定布局：

\begin{lstlisting}
block 0       boot/reserved
block 1       superblock
block 2..     redo log
随后          inode blocks
随后          free-block bitmap
随后          data blocks
\end{lstlisting}

superblock 字段定义在
\href{../../../kernel/fs.h}{\passthrough{\lstinline!kernel/fs.h!}}：

\begin{lstlisting}
magic, size, nblocks, ninodes,
nlog, logstart, inodestart, bmapstart
\end{lstlisting}

根目录 inode 号是 \passthrough{\lstinline!ROOTINO=1!}。每个 inode 有 12
个直接块和 1 个一级间接块， 最大文件块数为
\passthrough{\lstinline!12 + BSIZE/4 = 268!}，约 268
KiB。目录本身也是普通 inode， 数据是一组
\passthrough{\lstinline!\{inum, name[30]\}!} 目录项。

\section{运行时 I/O 层次}\label{ux8fd0ux884cux65f6-io-ux5c42ux6b21}

\begin{figure}[htbp]
\centering
\begin{tikzpicture}[node distance=6mm]
  \node[hbox] (sys) {系统调用\\open/read/write/unlink};
  \node[box,below=of sys] (name) {namei/nameiparent\\路径与目录项};
  \node[box,below=of name] (inode) {inode cache\\ilock/readi/writei/bmap};
  \node[box,below left=7mm and 13mm of inode] (log) {redo log\\begin\_op/log\_write/end\_op};
  \node[box,below right=7mm and 13mm of inode] (bc) {bcache\\bread/bwrite/brelse};
  \node[box,below=20mm of inode] (backend) {xv6fs block backend\\原始 SD / FAT32 FS.IMG / 内嵌 ramdisk};
  \draw[arr] (sys)--(name); \draw[arr] (name)--(inode); \draw[arr] (inode)--(log);
  \draw[arr] (inode)--(bc); \draw[arr] (log)--(bc); \draw[arr] (bc)--(backend);
\end{tikzpicture}
\caption{xv6fs 从 VFS/系统调用到块设备后端的 I/O 层次}
\label{fig:tiny-fs-io}
\end{figure}


\passthrough{\lstinline!fat32rw()!} 是当前 xv6fs
的兼容后端入口。它根据启动探测结果，把同一个 xv6fs 逻辑块映射到原始 SD
分区、FAT32 文件 \passthrough{\lstinline!FS.IMG!} 的 cluster
chain，或最终交给
\href{../../../kernel/drivers/block/ramdisk_embed.c}{\passthrough{\lstinline!ramdisk\_embed.c!}}
访问内嵌镜像。

\section{bcache}\label{bcache}

\href{../../../kernel/bio.c}{\passthrough{\lstinline!bio.c!}}
提供块缓存：

\begin{itemize}
\tightlist
\item
  key 是 \passthrough{\lstinline!(dev, blockno)!}；
\item
  cache miss 时复用 LRU 链表尾部且 \passthrough{\lstinline!refcnt==0!}
  的 buffer；
\item
  每个 buffer 用 sleeplock 保证同一块只有一个使用者修改；
\item
  \passthrough{\lstinline!bread()!} 在
  \passthrough{\lstinline!valid==0!} 时调用底层
  \passthrough{\lstinline!fat32rw(..., read)!}；
\item
  \passthrough{\lstinline!bwrite()!} 同步调用底层写；
\item
  日志中的脏块通过 \passthrough{\lstinline!bpin()!} 防止提交前被回收。
\end{itemize}

它缓存的是 1 KiB 文件系统块，不是文件级 page cache。使用 SD/FAT32
后端时，bcache 还会避免重复的介质访问；日志中的脏块通过 pin
保证提交前不会被替换。

\section{redo log}\label{redo-log}

\href{../../../kernel/log.c}{\passthrough{\lstinline!log.c!}} 实现物理块
redo log。文件系统修改调用：

\begin{lstlisting}
begin_op()
  -> 修改 inode/bitmap/directory/data 对应的 cached block
  -> log_write(buf)                 记录 home block，pin buffer
end_op()
  -> 最后一个 outstanding 操作触发 commit()
       1. write_log()               cache -> log data blocks
       2. write_head()              写 log header；事务提交点
       3. install_trans()           log blocks -> home blocks
       4. write_head(n=0)           清事务
\end{lstlisting}

启动时 \passthrough{\lstinline!initlog()!} 调用
\passthrough{\lstinline!recover\_from\_log()!}：若 log header 的
\passthrough{\lstinline!n>0!}，把日志块 重新安装到 home block，再清
header。

使用 SD 原始分区或 FAT32 \passthrough{\lstinline!FS.IMG!} 时，日志能让
xv6fs 修改跨重启恢复；使用内嵌 ramdisk
兜底时，重启会重新装入构建时镜像，因此运行时写入不会持久化。

\section{文件系统何时初始化}\label{ux6587ux4ef6ux7cfbux7edfux4f55ux65f6ux521dux59cbux5316}

\passthrough{\lstinline!x210\_main()!} 先完成
\passthrough{\lstinline!binit()!}、\passthrough{\lstinline!iinit()!} 和
\passthrough{\lstinline!fileinit()!}，然后创建
\passthrough{\lstinline!initproc!}。 第一次调度到普通用户进程时：

\begin{lstlisting}
forkret()
  -> fsinit(ROOTDEV)
       -> bread(dev, 1) 读取 superblock
       -> 校验 FSMAGIC
       -> initlog() / recover_from_log()
  -> usertrapret()
  -> initcode exec("/init")
\end{lstlisting}

\passthrough{\lstinline!fsinit()!} 放在进程上下文，是因为 bcache 的
sleeplock 和日志路径可能睡眠，不能在 尚无当前进程的早期
\passthrough{\lstinline!x210\_main()!} 中执行。

\section{\texorpdfstring{\texttt{/init}
与设备节点}{/init 与设备节点}}\label{init-ux4e0eux8bbeux5907ux8282ux70b9}

Tiny210 使用
\href{../../../kernel/arch/arm/init.c}{\passthrough{\lstinline!kernel/arch/arm/init.c!}}：

\begin{enumerate}
\def\labelenumi{\arabic{enumi}.}
\tightlist
\item
  创建 \passthrough{\lstinline!/dev!}；
\item
  创建 \passthrough{\lstinline!/dev/console!}（major 1）和
  \passthrough{\lstinline!/dev/ttyS0!}（major 3）；
\item
  \passthrough{\lstinline!tty\_attach(0)!}，打开 ttyS0 并复制为 fd
  0/1/2；
\item
  fork 子进程并 \passthrough{\lstinline!exec("/bin/sh")!}；
\item
  shell 退出后重新启动。
\end{enumerate}

\passthrough{\lstinline!mkdir/mknod!}
遇到镜像中已存在的对象时返回失败；当前 init 没有严格检查 mkdir 和 mknod
的返回值，因此已有正确节点不会阻止继续启动，但错误类型/major 也不会被
自动修复。更稳定的方向是在 \passthrough{\lstinline!mkfs!}
时创建固定目录和设备命名空间。

\section{后续改进}\label{ux540eux7eedux6539ux8fdb}

当前已完成 SDHCI、FAT32、VFS、只读 ext2 和 xv6fs 后端选择。下一步应：

\begin{enumerate}
\def\labelenumi{\arabic{enumi}.}
\tightlist
\item
  把 xv6fs 的兼容 \passthrough{\lstinline!fat32rw()!} 入口替换为标准
  block-device operations。
\item
  让 bcache 按设备逻辑块大小工作，减少 FAT32 file-sector 的重复转换。
\item
  完善 ext2 写入、缓存一致性和挂载生命周期。
\item
  加入可靠的 flush/barrier，保证真实介质上 journal commit 的持久化顺序。
\item
  增加断电恢复测试，区分 xv6fs 日志、FAT32 容器元数据和 SD
  卡内部缓存三层风险。
\end{enumerate}
